% =====================================================================
% Section 7: Discussion
% =====================================================================

\section{Discussion}
\label{sec:discussion}

This section interprets the empirical results
(\S\ref{subsec:what-mass-does}), positions the algorithm in the
landscape of network architectures
(\S\ref{subsec:protocol-implications}), records its limitations
honestly (\S\ref{subsec:limitations}), and outlines work the
authors believe is worth pursuing
(\S\ref{subsec:future-work}).

\subsection{What the mass mechanism actually does}
\label{subsec:what-mass-does}

The improvements documented in \S\ref{sec:results} are produced by
an algorithm that, stripped of motivational language, has a
straightforward operational description:
\textit{a packet that has already crossed congested edges becomes
internally more expensive to route, and the routing algorithm
systematically prefers cooler alternatives for it.}

This produces an effect that is \emph{not} reducible to either of
the two standard categories of routing improvement. It is not
load-aware in the classical sense (LASP, ECMP) because two packets
arriving at the same node from different histories receive
different routing recommendations. It is not learned in the
reinforcement-learning sense because no policy is trained or
updated; the algorithm is a fixed dynamic program parameterized by
$\kappa$ alone. What the mass mechanism does is closer in spirit
to a form of \emph{trajectory-conditioned scheduling}: routing
decisions for a packet depend on what that packet has already
experienced, not only on the current network state.

The empirical signature of this mechanism is consistent with the
interpretation: improvement scales with the amount of congestion
(\S\ref{subsec:vanishing}), is largest where the topology offers
genuine alternatives (real backbones, regular lattices), and
vanishes in regimes where the network has spare capacity
everywhere. The mechanism is also \emph{conservative}: average
hop counts increase by only $0.04$ on GÉANT, indicating that
the algorithm is not trading capacity for safety in any
substantial way. It is using existing capacity more
intelligently.

\subsection{Implications for protocol design}
\label{subsec:protocol-implications}

EMMET-DP is an admission-time, controller-level routing
algorithm. It assumes access to global topology and to per-edge
state, makes a single routing decision per packet \emph{at the
moment the packet is admitted to the network}, and emits a fixed
forwarding path. The mass scalar evolves internally during the
dynamic-programming search over candidate paths; once the path
is selected, packets do not reroute mid-flight under the
evaluation in this paper. This places it in the same architectural family
as software-defined networking control planes such as B4 and
SWAN~\cite{jain2013b4,hong2013swan}, and unlike data-plane schemes
such as ECMP that distribute routing decisions per-flow at line
rate.

For deployment, three architectural observations follow:

\paragraph{The mass scalar is cheap.}
The per-packet state is a single floating-point value (or, in our
implementation, a 5-bit bucket index). Carrying it in a header
field costs less than carrying a flow identifier, and updating it
requires a single fused-multiply-add per hop. There is no
fundamental obstacle to in-band carriage in protocols where
header extension is feasible (IPv6 hop-by-hop options,
SRv6~\cite{filsfils2018srv6}).

\paragraph{The DP is centralizable, not on the critical path.}
The dynamic-programming router runs at admission time, not at
forwarding time. In a controller-augmented network, the controller
emits a forwarding label encoding the precomputed path; the data
plane forwards along that label without re-running the DP. The
algorithm is therefore compatible with line-rate forwarding even
on hardware that cannot evaluate the DP itself.

\paragraph{Mass dynamics generalize beyond admission.}
A more ambitious deployment would have intermediate routers update
the mass scalar as the packet traverses, allowing reroute decisions
mid-flight in response to mass-bucket transitions. We have not
evaluated this regime; the present paper restricts itself to
admission-time routing, where the algorithm reduces cleanly to
existing controller-plane infrastructure.

\subsection{Limitations}
\label{subsec:limitations}

We record limitations of the present work in the order in which a
practitioner is likely to encounter them.

\paragraph{Synthetic workload.}
Our traffic model is uniformly random source--destination pairs at
fixed step intervals. Real datacenter and ISP traffic shows
heavy-tailed flow sizes, locality, and time-of-day patterns that
our experiments do not capture. The improvements we document
should therefore be read as bounds on \emph{algorithmic}
performance under controlled conditions, not as predicted
production gains.

\paragraph{Single-domain evaluation.}
Both real backbones we evaluate (GÉANT, Abilene) are
single-AS academic networks. Inter-AS routing introduces policy
constraints, BGP-driven path selection, and economic
considerations that EMMET-DP does not currently model. The
algorithm is single-domain by construction; multi-domain
extension is non-trivial and out of scope.

\paragraph{Static edge attributes.}
Edge latency and capacity are sampled once per topology
realization and held fixed. Real networks experience link flaps,
maintenance windows, and time-varying capacity (e.g., wireless,
shared queues). Extending the algorithm to handle attribute drift
is mechanically straightforward (the loss snapshot $\theta$
already responds to recent loss events) but requires an
evaluation we have not performed.

\paragraph{Hyperparameter discovery.}
The headline value $\kappa = 1.0$ was selected from a sweep over a
single discrete grid (\S\ref{subsec:kappa-sweep}). Production
deployment would benefit from automated $\kappa$ selection per
deployment, possibly tied to a measurable global quantity such as
$T(G)$. We do not provide such a procedure here.

\paragraph{Per-packet ordering effects.}
Our simulation processes traffic events sequentially within a
seed, so packet $i+1$ observes the network state left by packet
$i$. In a real network, packets are injected concurrently and
observe an aggregated state. The ordering effect is bounded by
the load-decay parameter $d_\lambda = 0.9$, but a more rigorous
treatment would model packet arrivals as a continuous-time
process. This is on the future-work list.

\subsection{Future work}
\label{subsec:future-work}

Three directions appear to us most promising:

\paragraph{Generalization to other domains.}
The mass-aware routing structure—per-agent state, multiplicative
update on encountering friction, dynamic programming over a
constrained state space—is not specific to packet networks. The
same template applies to robotics navigation through hazard
fields, logistics under congestion, and epidemic-aware routing in
contact graphs. We are preparing a companion paper that lifts the
algorithm out of the networking context and treats it as a
general framework for trajectory-aware planning under
medium-induced cost dynamics.

\paragraph{Online learning of $\kappa$.}
Rather than fixing $\kappa$ at deployment time, the controller
could adjust it from observed loss rates and global temperature,
using a slow outer loop while EMMET-DP runs the fast inner loop.
This sits naturally between the fully-fixed regime evaluated here
and the fully-learned regime of reinforcement-learning routers,
and may capture most of the benefits of both.

\paragraph{Hardware-accelerated DP.}
The dynamic program is regular: $|V| \cdot H \cdot B$ states with
constant work per state. This makes it a candidate for
acceleration on FPGAs, network processors, or eBPF in-kernel
paths. We have not benchmarked any of these. A latency floor
below 10 microseconds per route would put EMMET-DP within reach
of inline use rather than admission-time precomputation.

A final note: this work argues that giving packets a small amount
of internal state, and carrying that state into routing decisions,
buys real improvement on a stateless-baseline that already uses
all available network-level information. The mechanism is simple,
cheap, and physically motivated. Whether it generalizes beyond
the regimes we evaluate is, of course, an empirical question—but
the evidence here is, to our reading, encouraging enough to
warrant the larger evaluation we are planning.
